\documentclass[tikz, border=10pt]{standalone}
\usepackage{pgfplots}
\pgfplotsset{compat=1.18}
\usetikzlibrary{patterns} % Librería necesaria para el patrón de rayas

\begin{document}
	\definecolor{maincolor1}{RGB}{1,82,142}
	\definecolor{maincolor2}{RGB}{21,182,184}
	\definecolor{mygrey}{rgb}{0.41, 0.41, 0.41}
	\definecolor{discardred}{RGB}{210, 60, 60}
	
	\begin{tikzpicture}
		\begin{axis}[
			title={$f(x) = x^2 - 3.2x + 3$ \quad (It. 0)},
			xlabel={$x$},
			ylabel={$f(x)$},
			xmin=-0.2, xmax=3.2,
			ymin=-0.1, ymax=3.3,
			axis lines=left,
			grid style={line width=.1pt, draw=gray!20},
			major grid style={line width=.2pt, draw=gray!50},
			minor tick num=1,
			legend pos=north west,
			legend style={font=\small, fill=white, fill opacity=0.9, draw opacity=1},
			width=12cm,
			height=8cm
			]
			
			% 1. Intervalo inicial completo de incertidumbre [0.0, 3.0] (Fondo gris sin descartes)
			\addplot[
			draw=none, 
			fill=mygrey, 
			opacity=0.1
			] coordinates {
				(0.0, -0.1) (0.0, 3.3) (3.0, 3.3) (3.0, -0.1)
			};
			
			% 2. Curva continua f(x)
			\addplot [
			domain=0:3, 
			samples=100, 
			color=maincolor1, 
			thick
			] {x^2 - 3.2*x + 3};
			
			% 3. Líneas proyectadas para los límites iniciales a0 = 0.0 y b0 = 3.0
			\draw[gray!50, dashed] (axis cs:0.0, -0.1) -- (axis cs:0.0, 3.0);
			\draw[gray!50, dashed] (axis cs:3.0, -0.1) -- (axis cs:3.0, 2.4);
			
			% 4. Puntos evaluados en los extremos del intervalo inicial
			\addplot [
			only marks, 
			mark=*, 
			mark size=3pt, 
			color=maincolor2
			] coordinates {
				(0.0, 3.0)
				(3.0, 2.4)
			};
			
			% 5. Mínimo analítico exacto (1.6, 0.44)
			\addplot [
			only marks, 
			mark=diamond*, 
			mark size=4pt, 
			color=maincolor2,
			thick
			] coordinates {(1.6, 0.44)};
			
			% Etiquetas en el eje X para a0 y b0
			\node[below, font=\footnotesize, maincolor1] at (axis cs:0.0, -0.1) {$a_0 = 0$};
			\node[below, font=\footnotesize, maincolor1] at (axis cs:3.0, -0.1) {$b_0 = 3$};
			
		\end{axis}
	\end{tikzpicture}
\end{document}